% ============================================================
\section{Dados: corpus, estudo de viabilidade e amostra}\label{sec:dados}
% ============================================================

Os vértices das redes são ocorrências de tokens em parágrafos reais da Wikipédia em português.
Esta seção descreve como o corpus é montado (etapa \etapa{corpus}), o estudo de viabilidade que
dimensionou a amostra e o desenho final da amostra (etapa \etapa{sample}).

% ------------------------------------------------------------
\subsection{Corpus}\label{sec:corpus}

\paragraph{Por que a Wikipédia.} É texto real, em português, com licença aberta e
organizado por categorias, o que permite definir \emph{temas} de forma reproduzível. Cada
artigo tem um identificador de revisão (\cod{revid}); guardando esses identificadores, o mesmo
texto pode ser baixado de novo no futuro, mesmo que o artigo tenha mudado: o comando
\cod{corpus-rebuild} baixa por \cod{revid} as revisões dos artigos mantidos e as limpa com o
mesmo código e na mesma ordem da etapa \etapa{corpus}.

\paragraph{Temas.} São oito, cada um definido pela árvore de uma categoria raiz: Física,
Biologia, Economia, Política, Ciência da computação, Música, Futebol (tema ``esporte'') e
Geografia. O tema cumpre dois papéis: é um dos rótulos de contexto da P3 e serve como
\emph{indicador aproximado do sentido} das palavras-alvo na P4 (\emph{banco} num artigo de
computação tende a ser banco de dados; em economia, instituição financeira). Os temas foram
escolhidos distantes entre si, para que o vocabulário de cada um seja característico, e de
modo que cada palavra-alvo tenha sentidos diferentes em pelo menos dois deles.

\paragraph{Busca por categorias.} A partir de cada raiz, uma busca em largura percorre as
subcategorias até a profundidade 2 e junta até 600 títulos por tema. Só entram páginas do
domínio principal (artigos); subcategorias de manutenção ou de usuário (títulos que contêm
``Artigos'', ``Esboço'', ``Páginas'', ``Predefinições'', ``Wikipédia'', ``Usuário'', ``Listas''
ou ``Categoria:!'') não são seguidas, e títulos ``Lista de\ldots'' são descartados
(Seção~\ref{sec:prob-categorias}). Todas as respostas cruas da API ficam num cache em disco
(\cod{data/raw/wiki}), o que torna o download retomável e a limpeza refazível sem rede. No
modo offline (\cod{GENDER\_NETWORKS\_OFFLINE=1}), uma resposta que falta no cache interrompe a
etapa, em vez de baixar uma revisão mais nova.

\paragraph{Categorias extras para os sentidos fracos.} No estudo de viabilidade, alguns
sentidos das palavras-alvo apareciam muito pouco no tema que os representa: \emph{nota} em
Economia (7 ocorrências, no sentido de cédula), \emph{carga} em Economia (6, carga
tributária), \emph{órgão} em Biologia (6) e \emph{banco} em Geografia (4, banco de areia);
\emph{matriz} e \emph{planta} também tinham um segundo sentido fraco. Para esses casos, o
corpus inclui categorias extras ligadas a um dos oito temas (por exemplo, Numismática e Impostos
em Economia; Anatomia humana e Botânica em Biologia; Geomorfologia e Oceanografia em
Geografia; Bancos de dados e Matrizes em Computação), com busca até a profundidade 1 e até 60
artigos por categoria. A lista completa está em \cod{configs/experiment.yaml}.

\paragraph{Páginas descartadas.} Um artigo é descartado quando:
\begin{itemize}
    \item é alcançado a partir de \textbf{mais de um tema} (contando as categorias extras). A
    busca registra todos os temas que chegam a cada título. Uma página multitema (por exemplo,
    Economia política) misturaria o vocabulário de dois temas e borraria justamente o rótulo
    usado na P3 e na P4;
    \item é uma \textbf{biografia}, detectada pelas predefinições de infobox
    (\cod{Info/Biografia}, \cod{Info/Futebolista}, \cod{Info/Político}, \cod{Info/Músico} e
    outras). Biografias aparecem em quase todas as árvores de categorias e falam de datas,
    lugares e carreiras, um vocabulário que não é do tema;
    \item é uma página de \textbf{desambiguação}, não tem texto, repete um artigo já escolhido
    ou não sobra nenhum parágrafo depois da limpeza.
\end{itemize}
Os artigos restantes de cada tema são embaralhados com semente fixa e baixados em lotes de 50
até completar a cota de 150 por tema (60 por categoria extra).

\paragraph{Limpeza.} As seções de material de apoio do \emph{wikitext} (Referências, Ligações
externas, Ver também, Bibliografia, Notas, Fontes e variantes) e suas subseções são removidas;
seções de conteúdo que vêm depois delas são mantidas (em artigos de música, por exemplo,
``Notas'' pode vir antes de ``História''). O texto é então convertido em texto com o \cod{mwparserfromhell}. Antes da conversão, a árvore sintática é
reescrita em três pontos, porque a conversão padrão (\cod{strip\_code}) mantém o conteúdo de
marcas que não são texto corrido:
\begin{itemize}
    \item marcas de bloco (\cod{<ref>}, \cod{<gallery>}, \cod{<timeline>}, \cod{<pre>},
    tabelas e outras) e links para arquivos e categorias são removidos sem deixar rastro; sem
    isso, o texto das referências entrava nos parágrafos (Seção~\ref{sec:prob-ref});
    \item marcas que não podem virar palavras (\cod{<math>}, \cod{<chem>}, \cod{<score>}) e itens
    de lista fazem a linha inteira ser descartada: apagar uma fórmula deixaria um buraco na
    frase (``a energia é dada por , onde\ldots''); fórmulas triviais (um símbolo ou um número,
    como \cod{<math>v</math>}) são mantidas como texto;
    \item predefinições em linha cujo texto visível é um argumento (\cod{formatnum},
    \cod{nowrap}, conversões de unidades como \cod{converter}) são trocadas por esse texto; as
    demais predefinições (infoboxes, navegação, pronúncia, datas formatadas) são descartadas, e
    os restos que deixam (``( )'', ``, ,'') são limpos por expressões regulares.
\end{itemize}

\paragraph{Filtros de parágrafo.} Fica só a prosa corrida: parágrafos com pelo menos 40
palavras, que terminam como frase (``.'', ``!'' ou ``?'', eventualmente seguidos de aspas ou
parênteses), não começam com marcação de tabela ou lista e têm no máximo 8\% de dígitos
(o que elimina tabelas de resultados e cronologias). Parágrafos repetidos são removidos pelos
primeiros 80 caracteres em minúsculas (textos copiados entre artigos são comuns). Cada
parágrafo é dividido em sentenças por expressão regular, com uma lista de abreviações do
português (``Sr.'', ``séc.'', ``p.ex.'', iniciais de nomes\ldots), para dar o rótulo de
sentença usado na P3.

\paragraph{Saídas.} \cod{data/\allowbreak corpus/\allowbreak articles.jsonl} (todos os candidatos, com o motivo de
descarte), \cod{data/\allowbreak corpus/\allowbreak paragraphs.jsonl} (parágrafos mantidos, com os limites das
sentenças) e \cod{data/\allowbreak corpus/\allowbreak manifest.csv}, a lista versionada de \cod{pageid},
\cod{revid}, título, tema e categoria de origem.

\subsubsection*{Estatísticas do corpus}

O corpus final tem \nArtigos{} artigos (de \nArtigosCandidatos{} candidatos), \nParagrafos{}
parágrafos, \nPalavrasCorpus{} palavras, \nTokensCorpus{} tokens do Qwen3 e \nTiposCorpus{}
tipos distintos.

\tabelaopcional[etapa \etapa{report} a partir de \etapa{corpus}]{corpus_temas}
    {Artigos, parágrafos e palavras por tema, e artigos descartados por motivo.}
    {tab:corpus-temas}
\figuraopcional[etapa \etapa{report} a partir de \etapa{corpus}]{corpus_paragrafos}
    {Parágrafos por tema e distribuição do tamanho dos parágrafos (em palavras).}
    {fig:corpus-paragrafos}
\figuraopcional[etapa \etapa{report} a partir de \etapa{corpus}]{corpus_descartes}
    {Páginas candidatas descartadas por motivo (multitema, biografia, desambiguação, sem
    parágrafos, duplicada), por tema.}
    {fig:corpus-descartes}

% ------------------------------------------------------------
\subsection{Estudo de viabilidade}\label{sec:viabilidade}

Antes de implementar o pipeline, um estudo rápido (scripts em \cod{scripts/feasibility/})
respondeu a três perguntas: (i) a meta de cerca de 15 mil vértices é alcançável com parágrafos
inteiros? (ii) Parágrafos inteiros deixam palavras repetidas suficientes, em temas diferentes,
para a P4? (iii) Que tokenizer, e portanto que modelo, usar?

\paragraph{Dados do estudo.} O \cod{fetch\_wiki.py} baixou uma versão simplificada do corpus:
os mesmos oito temas, busca até a profundidade 2, 120 artigos por tema, limpeza só com
\cod{strip\_code} e os mesmos filtros de parágrafo. Resultaram \textbf{958 artigos} e
\textbf{7.876 parágrafos} (Biologia 965, Computação 1.034, Economia 1.147, Esporte 836, Física
948, Geografia 785, Música 1.032, Política 1.129), com \textbf{1,28 milhão de tokens} do
Qwen3 (1.276.475) e cerca de 23 mil tipos distintos. Três tokenizers foram comparados: o do
Qwen3-4B (idêntico ao do Qwen3-4B-Base, 151.669 ids), o do Tucano-2b4 (treinado em português) e
o do GPT-2 small português (o modelo do piloto).

\paragraph{Amostragem natural.} O \cod{simulate\_sample.py} simula a amostragem mais simples:
parágrafos inteiros, em rodízio entre os temas, com todos os tokens como vértices e um limite
$f_{\max}$ por tipo, até chegar a 15 mil vértices. Também simula uma variante
``direcionada'', que começa pelos parágrafos que contêm as palavras-alvo (até 12 por palavra e
tema) e depois segue a ordem natural. As saídas estão em \cod{sim\_<modelo>.txt}.

\paragraph{Desenho híbrido.} O \cod{sim\_hybrid.py} simula a alternativa adotada: um
\emph{núcleo} denso de parágrafos inteiros mais estratos esparsos de ocorrências escolhidas
(palavras-alvo, controles e palavras de conteúdo presentes em vários temas), em que só as
posições sorteadas viram vértices e o resto do parágrafo é contexto. As saídas estão em
\cod{hyb\_<modelo>.txt}.

Nas contagens a seguir, \emph{palavra de conteúdo} é um tipo cuja categoria mais comum é
palavra inteira, só com letras, com pelo menos 3 letras e fora de uma lista de palavras
funcionais. O critério da P4 é ter $f_t\ge10$ na amostra e pelo menos 3 ocorrências em pelo
menos 3 temas.

\begin{table}[htbp]
\centering
\caption{Estudo de viabilidade: tokenizers e amostragem natural até 15 mil vértices com
$f_{\max}=50$. Fontes: linha \cod{pool:} e bloco \cod{f\_max=50, amostragem natural} de
\cod{scripts/feasibility/sim\_<modelo>.txt}.}\label{tab:viab-natural}
\small
\begin{tabular}{@{}l r r r@{}}
\toprule
& \textbf{Qwen3-4B} & \textbf{Tucano-2b4} & \textbf{GPT-2 small PT} \\
\midrule
Tokens no \emph{pool} de 7.876 parágrafos & 1.276.475 & 953.290 & 939.895 \\
Tokens por palavra & 1,78 & 1,33 & 1,31 \\
Palavras-alvo candidatas que são um token só & 12 de 16 & 16 de 16 & 16 de 16 \\
\midrule
Parágrafos para chegar a 15 mil vértices & 101 & 134 & 160 \\
Tipos distintos na amostra & 3.693 & 6.144 & 6.213 \\
Vértices com $f_t\le2$ & 19\% & 40\% & 39\% \\
Vértices que são palavras inteiras & 28\% & 54\% & 72\% \\
Tipos com $f_t\ge10$ & 293 & 183 & 170 \\
Conteúdo com $f_t\ge10$ e $\ge3$ oc.\ em $\ge3$ temas & 4 & 5 & 9 \\
Alvos com $\ge10$ oc.\ em $\ge2$ temas & 0 & 0 & 0 \\
\midrule
\multicolumn{4}{@{}l}{\emph{Variante direcionada (alvos primeiro), $f_{\max}=50$}} \\
Conteúdo com $f_t\ge10$ e $\ge3$ oc.\ em $\ge3$ temas & 10 & 10 & 21 \\
Alvos com $\ge10$ oc.\ em $\ge2$ temas & 2 & 1 & 2 \\
\bottomrule
\end{tabular}
\end{table}

\figuraopcional[\cod{scripts/feasibility/plot\_feasibility.py}]{viabilidade_tokenizers}
    {Comparação dos tokenizers no \emph{pool} do estudo de viabilidade: tokens por palavra e
    fração dos vértices que são palavras inteiras na amostragem natural ($f_{\max}=50$). O
    Qwen3 (em destaque) é o tokenizer do modelo escolhido. Fonte:
    \cod{scripts/feasibility/sim\_<modelo>.txt}.}
    {fig:viab-tokenizers}

A Tabela~\ref{tab:viab-natural} e a Figura~\ref{fig:viab-tokenizers} mostram três coisas.

\textbf{A meta de vértices é fácil.} Com qualquer tokenizer, 15 mil vértices saem de 100 a 160
parágrafos (de 18 a 21 mil tokens), uma fração pequena do \emph{pool}. Com $f_{\max}=20$ bastam
de 127 a 177 parágrafos.

\textbf{Parágrafos inteiros não servem para a P4.} Mesmo com 15 mil vértices, só de 4 a 9
palavras de conteúdo têm $f_t\ge10$ com ocorrências espalhadas por 3 temas, e \emph{nenhuma}
palavra-alvo chega a 10 ocorrências em 2 temas. O motivo é a lei de Zipf: num trecho pequeno de
texto, só palavras funcionais e pontuação se repetem muito; palavras de conteúdo aparecem uma
ou duas vezes. Começar pelos parágrafos que contêm as palavras-alvo (variante direcionada)
melhora pouco (de 10 a 21 palavras de conteúdo, 1 ou 2 alvos).

\textbf{O Qwen3 fragmenta muito o português.} São 1,78 tokens por palavra, contra 1,33 do
Tucano e 1,31 do GPT-2 português, e só 28\% dos vértices são palavras inteiras (54\% e 72\% nos
outros). Quatro das 16 palavras-alvo candidatas não são um token só no Qwen3 (\emph{corrente},
\emph{onda}, \emph{célula} e \emph{núcleo}; por exemplo, ``\esp célula'' vira
\tok{\esp cél}${}+{}$\tok{ula}) e saíram da lista, que ficou com 12. Essa foi a principal
desvantagem pesada contra o Qwen3 na escolha do modelo (Seção~\ref{sec:mod-escolha}).

\begin{table}[htbp]
\centering
\caption{Estudo de viabilidade: desenho híbrido com 15 mil vértices e $f_{\max}=50$
(núcleo de parágrafos inteiros, 12 ou 16 alvos, 6 controles e 150 palavras multitema com 20
ocorrências cada). Fonte: \cod{scripts/feasibility/hyb\_<modelo>.txt}.}\label{tab:viab-hibrido}
\small
\begin{tabular}{@{}l r r r@{}}
\toprule
& \textbf{Qwen3-4B} & \textbf{Tucano-2b4} & \textbf{GPT-2 small PT} \\
\midrule
Vértices & 15.104 & 15.002 & 15.099 \\
\quad núcleo & 11.251 & 10.954 & 11.045 \\
\quad alvos & 553 & 748 & 754 \\
\quad controles & 300 & 300 & 300 \\
\quad multitema & 3.000 & 3.000 & 3.000 \\
Parágrafos inteiros no núcleo & 78 & 93 & 107 \\
Vértices por parágrafo do núcleo (média) & 144 & 118 & 103 \\
Parágrafos tocados no total & 2.894 & 2.943 & 3.017 \\
Tokens a processar no modelo & 329.555 & 243.959 & 248.579 \\
Tipos distintos & 3.257 & 4.904 & 4.856 \\
Vértices com $f_t\le2$ & 17\% & 32\% & 30\% \\
Tipos com $f_t\ge10$ & 354 & 274 & 286 \\
Conteúdo com $f_t\ge10$ e $\ge3$ oc.\ em $\ge3$ temas & 169 & 172 & 173 \\
Alvos com $\ge10$ oc.\ em $\ge2$ temas & 3 de 12 & 7 de 16 & 7 de 16 \\
Alvos com $\ge5$ oc.\ em $\ge3$ temas & 11 de 12 & 14 de 16 & 14 de 16 \\
\bottomrule
\end{tabular}
\end{table}

\figuraopcional[\cod{scripts/feasibility/plot\_feasibility.py}]{viabilidade_amostragem}
    {Amostragem natural (parágrafos inteiros, $f_{\max}=50$) contra o desenho híbrido, nas duas
    contagens que decidem a P4: palavras de conteúdo repetidas em vários temas e palavras-alvo
    com ocorrências suficientes em dois temas. Fontes:
    \cod{scripts/feasibility/sim\_<modelo>.txt} e \cod{hyb\_<modelo>.txt}.}
    {fig:viab-amostragem}

\textbf{O desenho híbrido resolve.} Com o mesmo número de vértices, o núcleo continua com cerca
de 11 mil vértices em 78 a 107 parágrafos inteiros (o bastante para a P3 e para as métricas
globais), e o número de palavras de conteúdo com $f_t\ge10$ espalhadas por 3 temas sobe de
4--9 para cerca de 170 (Tabela~\ref{tab:viab-hibrido} e Figura~\ref{fig:viab-amostragem}). O
custo é processar mais texto no modelo: cerca de 330 mil tokens com o Qwen3, porque cada
ocorrência esparsa exige o parágrafo até a posição dela.

Na simulação, os alvos ainda ficaram fracos (3 de 12 com $\ge10$ ocorrências em 2 temas), mas
por uma razão corrigível: as 50 ocorrências de cada alvo foram sorteadas em rodízio entre os
\emph{oito} temas, e não concentradas nos temas em que a palavra tem sentidos diferentes. A
Tabela~\ref{tab:viab-pool} mostra que o \emph{pool} tem ocorrências suficientes para cotas por
\emph{tema de sentido} na maioria dos casos, e onde faltam (\emph{banco} em Geografia,
\emph{órgão} em Biologia, \emph{nota} e \emph{carga} em Economia, \emph{planta} em Economia,
\emph{matriz} em Computação), as categorias extras da Seção~\ref{sec:corpus} entram. A
simulação também tinha dois defeitos de desenho, corrigidos no desenho final
(Seção~\ref{sec:prob-simulacao}): o limite por tipo do núcleo era aplicado por ordem de
chegada, e cada parágrafo do núcleo vinha de um artigo diferente.

\begin{table}[htbp]
\centering
\caption{Ocorrências das palavras-alvo e de controle no \emph{pool} inteiro do estudo de
viabilidade (tokenizer do Qwen3, sem limite por tipo), por tema. Em negrito, os temas de
sentido usados nas cotas do desenho final (Tabela~\ref{tab:cotas}). Fonte: bloco \cod{pool
inteiro} de \cod{scripts/feasibility/sim\_Qwen3-4B.txt}.}\label{tab:viab-pool}
\footnotesize
\setlength{\tabcolsep}{4.5pt}
\begin{tabular}{@{}l r r r r r r r r r@{}}
\toprule
& \textbf{Total} & \textbf{Bio} & \textbf{Comp} & \textbf{Econ} & \textbf{Esp} & \textbf{Fís}
& \textbf{Geo} & \textbf{Mús} & \textbf{Pol} \\
\midrule
banco & 66 & 9 & \textbf{19} & \textbf{24} & 6 & 0 & \textbf{4} & 3 & 1 \\
campo & 335 & 34 & 26 & 22 & \textbf{74} & \textbf{82} & \textbf{72} & 11 & 14 \\
órgão & 57 & \textbf{6} & 1 & 6 & 9 & 0 & 0 & \textbf{25} & \textbf{10} \\
estado & 282 & 24 & 39 & 28 & 13 & \textbf{84} & 19 & 20 & \textbf{55} \\
carga & 53 & 9 & 4 & \textbf{6} & 0 & \textbf{22} & 0 & 3 & 9 \\
rede & 133 & 4 & \textbf{71} & 15 & 2 & 11 & \textbf{15} & 3 & 12 \\
nota & 41 & 2 & 5 & \textbf{7} & 0 & 2 & 0 & \textbf{24} & 1 \\
capital & 271 & 1 & 0 & \textbf{201} & 18 & 1 & 9 & 6 & \textbf{35} \\
família & 70 & \textbf{10} & 10 & 4 & 4 & 13 & 2 & 13 & \textbf{14} \\
massa & 140 & 2 & 3 & 10 & 3 & \textbf{59} & \textbf{39} & 14 & 10 \\
matriz & 27 & \textbf{10} & \textbf{3} & 2 & 0 & 5 & 0 & 6 & 1 \\
planta & 35 & \textbf{25} & 0 & \textbf{0} & 0 & 0 & 9 & 0 & 1 \\
\midrule
ano & 310 & 13 & 15 & \textbf{36} & \textbf{82} & 35 & 30 & \textbf{65} & 34 \\
século & 356 & 29 & 12 & 31 & 11 & \textbf{86} & \textbf{66} & 59 & \textbf{62} \\
grande & 527 & \textbf{58} & 48 & 67 & 66 & \textbf{88} & 49 & 77 & 74 \\
primeiro & 510 & 24 & \textbf{52} & 39 & \textbf{134} & 63 & 48 & 92 & 58 \\
importante & 247 & \textbf{53} & 20 & \textbf{47} & 9 & 24 & 23 & 36 & 35 \\
água & 249 & \textbf{102} & 0 & 24 & 2 & 47 & \textbf{65} & 1 & 8 \\
\bottomrule
\end{tabular}
\end{table}

Os números do estudo são aproximados por construção: o corpus do estudo não removia páginas
multitema nem biografias, a limpeza ainda deixava restos de referências e as contagens
ignoram os limites por parágrafo e por artigo do desenho final. Eles serviram para escolher o
desenho e o modelo; a composição real da amostra está na Seção~\ref{sec:amostra}.

% ------------------------------------------------------------
\subsection{Amostra híbrida}\label{sec:amostra}

A amostra final tem cerca de 15 mil vértices, $f_{\max}=50$ e sementes fixas (semente 438),
com quatro estratos sorteados em ordem por um único gerador aleatório.

\paragraph{1. Núcleo (P1--P3).} Por tema, 5 artigos $\times$ 2 parágrafos (cerca de 80
parágrafos e 11 mil vértices), com todos os tokens que não são só espaço como vértices. Dois
parágrafos por artigo fazem \emph{artigo} e \emph{parágrafo} serem rótulos distintos na P3 (na
simulação eles quase coincidiam). O limite $f_{\max}$ é aplicado por \textbf{sorteio uniforme}
entre as ocorrências de cada tipo acima do limite, depois de juntar todo o núcleo, e não por
ordem de chegada: assim as ocorrências mantidas de ``de'' ou ``,'' se espalham por todos os
parágrafos, em vez de serem as dos primeiros parágrafos. O núcleo é o único estrato em que um
parágrafo aparece inteiro, e por isso a P3 é analisada só nele.

\paragraph{2. Palavras-alvo (P4).} Doze palavras com sentidos diferentes em temas diferentes:
\emph{banco}, \emph{campo}, \emph{órgão}, \emph{estado}, \emph{carga}, \emph{rede},
\emph{nota}, \emph{capital}, \emph{família}, \emph{massa}, \emph{matriz} e \emph{planta}. Cada
uma tem cotas por \textbf{tema de sentido}: 2 ou 3 temas, com 16 a 25 ocorrências em cada
(Tabela~\ref{tab:cotas}). A cota conta as ocorrências do alvo que já estão no núcleo daquele
tema; só a diferença é sorteada em parágrafos fora do núcleo, com no máximo 1 ocorrência por
(tipo, parágrafo) e 3 por (tipo, artigo), para que um único artigo não domine a amostra de uma
palavra. As ocorrências que ficam no núcleo mantêm o estrato \cod{core}, mas recebem
\cod{target\_word} e \cod{sense\_theme} como as sorteadas; por isso as contagens por estrato
(\nVerticesAlvo{} alvos, \nVerticesControle{} controles) ficam abaixo da soma das cotas, sem que
falte ocorrência. Só conta o id da forma minúscula precedida de espaço
(Seção~\ref{sec:categorias}). Um alvo que não chegar a 10 ocorrências (\cod{min\_per\_sense})
em pelo menos 2 temas sai da lista, e as faltas por (alvo, tema) ficam registradas.

\paragraph{3. Controles (P4).} Seis palavras de sentido estável (\emph{ano}, \emph{século},
\emph{grande}, \emph{primeiro}, \emph{importante} e \emph{água}) com a \textbf{mesma
alocação} dos alvos: duas com três temas e cotas 17/17/16, quatro com dois temas e cotas
25/25. Se os controles tivessem outro número de temas ou de ocorrências, uma diferença entre
alvos e controles na P4 poderia vir só do desenho.

\paragraph{4. Multitema (P4).} 150 palavras de conteúdo com pelo menos 5 ocorrências em pelo
menos 3 temas fora do núcleo, completadas até 20 ocorrências no total (contando as que já
estão na amostra, no núcleo ou nos estratos anteriores), sorteadas em rodízio entre os temas.
Por isso o estrato multitema (\nVerticesMultitema{} vértices) fica abaixo de $150\times20$.
Elas ampliam a P4 para além das 18 palavras escolhidas à mão.

\begin{table}[htbp]
\centering
\caption{Cotas por tema de sentido do desenho final (\cod{configs/experiment.yaml}). São
parâmetros do desenho; as quantidades efetivamente obtidas estão na
Tabela~\ref{tab:amostra-cotas}.}\label{tab:cotas}
\small
\begin{tabular}{@{}l l@{\hspace{2.5em}}l l@{}}
\toprule
\textbf{Alvo} & \textbf{Cotas} & \textbf{Controle} & \textbf{Cotas} \\
\midrule
banco & Economia 17, Computação 17, Geografia 16 & ano & Economia 17, Esporte 17, Música 16 \\
campo & Física 17, Esporte 17, Geografia 16 & século & Física 17, Geografia 17, Política 16 \\
órgão & Biologia 17, Música 17, Política 16 & grande & Física 25, Biologia 25 \\
estado & Física 25, Política 25 & primeiro & Esporte 25, Computação 25 \\
carga & Física 25, Economia 25 & importante & Biologia 25, Economia 25 \\
rede & Computação 25, Geografia 25 & água & Biologia 25, Geografia 25 \\
nota & Música 25, Economia 25 & & \\
capital & Economia 25, Política 25 & & \\
família & Biologia 25, Política 25 & & \\
massa & Física 25, Geografia 25 & & \\
matriz & Biologia 25, Computação 25 & & \\
planta & Biologia 25, Economia 25 & & \\
\bottomrule
\end{tabular}
\end{table}

\paragraph{Sequências.} Cada parágrafo tocado pela amostra vira uma sequência
\tok{<|endoftext|>} $+$ tokens do parágrafo, truncada logo depois da última posição escolhida
nele (Seção~\ref{sec:mod-sequencias}). Nos estratos esparsos, os tokens antes da posição
sorteada são só contexto. A tabela \cod{occurrences.csv} tem uma linha por vértice, ordenada
por (sequência, posição), com origem (tema, artigo, revisão, parágrafo, sentença), posição e
faixa de posição, grupo de prefixo, token (id, texto, \emph{offsets}, vizinhos na sequência),
palavra e categoria, frequências e faixas na amostra e no corpus, estrato, palavra-alvo, tema
de sentido e contexto local. A ordem das linhas define o desempate por posição
(Seção~\ref{sec:red-empates}). Para a separação por sentido na P4, cerca de 20 ocorrências por
alvo são anotadas à mão em \cod{senses.csv}.

\paragraph{A amostra obtida.} Na última execução, a amostra tem \nVertices{} vértices em
\nSequencias{} sequências (\nTokensProcessados{} tokens a processar no modelo) e \nTipos{}
tipos distintos: \nVerticesNucleo{} no núcleo (\nParagrafosNucleo{} parágrafos),
\nVerticesAlvo{} de alvos, \nVerticesControle{} de controles e \nVerticesMultitema{}
multitema. Ficaram \nAlvosMantidos{} alvos na análise.

\tabelaopcional[etapa \etapa{report} a partir de \etapa{sample}]{amostra_composicao}
    {Composição da amostra por estrato e tema (vértices, parágrafos e tipos).}
    {tab:amostra-composicao}
\tabelaopcional[etapa \etapa{report} a partir de \etapa{sample}]{amostra_faixas}
    {Vértices por faixa de frequência (na amostra e no corpus) e por categoria de token.}
    {tab:amostra-faixas}
\tabelaopcional[etapa \etapa{report} a partir de \etapa{sample}]{amostra_cotas}
    {Ocorrências por palavra-alvo e controle em cada tema de sentido: cota, ocorrências que já
    estavam no núcleo, sorteadas fora dele e total obtido; alvos descartados por falta de
    ocorrências.}
    {tab:amostra-cotas}
\figuraopcional[etapa \etapa{report} a partir de \etapa{sample}]{amostra_composicao}
    {Composição da amostra: vértices por estrato, tema, faixa de frequência e categoria de
    token.}
    {fig:amostra-composicao}
